\section{Motivating analysis: covariance structure of inputs and cotangents}
\label{sec:motivation}

Before choosing which factors of \eqref{eq:family} to form or approximate, we
ask a prior empirical question on the baseline \muont{} trajectory itself:
\emph{what do the input second moments $C_t=\mathbb E[xx^{\Tr}]$ and the unused
output-cotangent second moments $B_t=\mathbb E[\delta\delta^{\Tr}]$ actually
look like across a transformer during training, and how fast do their
eigenspaces move?}

In a standard Llama-style block with seven linear projections
($q,k,v,o,\text{gate},\text{up},\text{down}$), shared consumed activations yield
four input moments per layer (\texttt{in:attn} for $q,k,v$, \texttt{in:o}, and
\texttt{in:mlp} for $\text{gate},\text{up}$ at hidden width $d=512$;
\texttt{in:down} post-SwiGLU at expansion width $d=2048$, for $128$ matrices
across $32$ layers), whereas distinct backward paths yield seven output moments
per layer (\texttt{out:q}, \texttt{out:k}, \texttt{out:v}, \texttt{out:o}, and
\texttt{out:down} at $d=512$; \texttt{out:gate} and \texttt{out:up} pre-SwiGLU
at $d=2048$, for $224$ matrices). Setting $L=\hat L=I$ while running the pooled
capture and cached-root machinery (\nullall{}, detailed in
Section~\ref{sec:method-controls}) executes pure input-side \muont{} updates
while recording all $352$ FP32 EMA moments ($\beta=0.938$, $N_t=36{,}608$
sampled tokens per update) and their cached eigenbases every $250$ updates
across the full $2.68\text{B}$-token ($5{,}120$-update) trajectory, completely
unconfounded by output preconditioning.

\subsection{Both sides are ill-conditioned, but their anisotropy lives in different coordinates}
\label{sec:mot-aniso}

For each symmetric PSD moment $M\in\mathbb R^{d\times d}$ across all $32$
layers, let $\tilde M = (d/\tr M)\,M$ denote its unit-mean-diagonal
normalization ($\frac{1}{d}\sum_{i=1}^d \tilde M_{ii}=1$) with eigenvalues
$\lambda_1,\dots,\lambda_d \ge 0$, and let
$\mathrm{Corr}_{ij} = M_{ij}/\sqrt{M_{ii}M_{jj}}$ denote its correlation matrix.
To separate spectral ill-conditioning from coordinate alignment, we evaluate six
complementary quantities:
\begin{enumerate}
\item \textbf{Effective rank share:}
  $\mathrm{erank}(M)/d = \frac{1}{d}\exp\!\bigl(-\sum_{i=1}^d p_i\log p_i\bigr)$
  with $p_i = \lambda_i / \sum_j \lambda_j$, measuring overall spectral
  concentration ($100\%$ for $I_d$).
\item \textbf{Coordinate scale spread:}
  $\mathrm{CV}_{\diag}(M) = \sqrt{\frac{1}{d}\sum_{i=1}^d (\tilde M_{ii}-1)^2}$,
  the coefficient of variation of per-neuron second moments.
\item \textbf{Off-diagonal correlation RMS:}
  $\rho_{\mathrm{rms}}(M) = \sqrt{\frac{1}{d(d-1)}\sum_{i\neq j}\mathrm{Corr}_{ij}^2}$,
  benchmarked against the single-update finite-sample correlation floor
  $N_t^{-1/2} = \motNoiseFloorCorrRms{}$ at $N_t=36{,}608$.
\item \textbf{Raw diagonal energy share:}
  $\mathrm{DiagE}(M) = \sum_{i=1}^d M_{ii}^2 / \sum_{i,j} M_{ij}^2$.
\item \textbf{Centered anisotropic diagonal share:} because even an isotropic
  matrix $\tilde M = I_d$ has $\mathrm{DiagE}(I_d)=100\%$ despite zero
  anisotropy, we subtract $I_d$ and invoke the exact Pythagorean decomposition
  $\|\tilde M - I_d\|_F^2 = \|\diag(\tilde M)-\mathbf 1\|_2^2 + \sum_{i\neq j}\tilde M_{ij}^2$
  to define
  \begin{equation}
  \label{eq:diaganiso}
  \mathrm{DiagAniso}(\tilde M)
  = \frac{\sum_{i=1}^d (\tilde M_{ii}-1)^2}{\sum_{i=1}^d (\tilde M_{ii}-1)^2 + \sum_{i\neq j}\tilde M_{ij}^2},
  \end{equation}
  the exact fraction of total departure from isotropy carried by coordinate
  variances rather than off-diagonal coupling.
\item \textbf{Diagonal root whitening fidelity:} with $A = \tilde M + \epsilon I_d$
  ($\epsilon = 10^{-3}$), $L = A^{-1/2}$, and
  $L_{\diag} = \Diag((\diag(\tilde M)+\epsilon\mathbf 1)^{-1/2})$,
  $R^2(A^{-1/2}) = 1 - \|L - L_{\diag}\|_F^2 / \|L - I_d\|_F^2$ measures how
  much of the spectral inverse square root's action is captured by elementwise
  diagonal scaling.
\end{enumerate}

\begin{table}[t]
\centering
\caption{Motivating covariance and subspace-drift audit of baseline input-side
\muont{} (\nullall{}, $L=\hat L=I$, $352$ FP32 EMA matrices across all $32$
layers). Columns $1$--$6$ report $32$-layer medians at mid-training (update
$2{,}500$, $1.31\text{B}$ tokens) so boundary layers $\text{L0}$--$\text{L3}$ do
not distort bulk summaries; $32$-layer means and $8$-checkpoint trajectory
averages ($0.13\text{B}$--$2.68\text{B}$ tokens) are reported in the text and
plotted in Figure~\ref{fig:motivation}. The final column ($\mathrm{Drift}_{16}$)
is the $32$-layer mean $250$-update top-$16$ Grassmann subspace drift between
updates $2{,}250$ and $2{,}500$ (Section~\ref{sec:mot-drift}).}
\label{tab:motivation}
\footnotesize\setlength{\tabcolsep}{4.8pt}
\resizebox{\textwidth}{!}{% Generated by scripts/analyze_mechanism.py (3-trajectory motivating audit: Muon, MuonT, WT-MLP-Lag). Do not edit.
\begin{tabular}{@{}l@{\hspace{6pt}}r@{\hspace{9pt}}rrr@{\hspace{9pt}}rrr@{\hspace{9pt}}rrr@{\hspace{9pt}}rrr@{\hspace{9pt}}rrr@{\hspace{9pt}}rrr@{\hspace{9pt}}rrr@{}}
\toprule
 & & \multicolumn{3}{c@{\hspace{9pt}}}{$\mathrm{erank}/d$ (\%)} & \multicolumn{3}{c@{\hspace{9pt}}}{$\mathrm{CV}_{\diag}$} & \multicolumn{3}{c@{\hspace{9pt}}}{$\rho_{\mathrm{rms}}$} & \multicolumn{3}{c@{\hspace{9pt}}}{$\mathrm{DiagE}$ (\%)} & \multicolumn{3}{c@{\hspace{9pt}}}{$\mathrm{DiagAniso}$ (\%)} & \multicolumn{3}{c@{\hspace{9pt}}}{$R^2(A^{-1/2})$ (\%)} & \multicolumn{3}{c@{}}{$\mathrm{Drift}_{16}$} \\
\cmidrule(r{9pt}){3-5}\cmidrule(r{9pt}){6-8}\cmidrule(r{9pt}){9-11}\cmidrule(r{9pt}){12-14}\cmidrule(r{9pt}){15-17}\cmidrule(r{9pt}){18-20}\cmidrule{21-23}
Moment & $d$ & \multicolumn{1}{c}{\tiny\taxid{Muon}} & \multicolumn{1}{c}{\tiny\taxid{MuonT}} & \multicolumn{1}{c@{\hspace{9pt}}}{\tiny\taxid{WT}} & \multicolumn{1}{c}{\tiny\taxid{Muon}} & \multicolumn{1}{c}{\tiny\taxid{MuonT}} & \multicolumn{1}{c@{\hspace{9pt}}}{\tiny\taxid{WT}} & \multicolumn{1}{c}{\tiny\taxid{Muon}} & \multicolumn{1}{c}{\tiny\taxid{MuonT}} & \multicolumn{1}{c@{\hspace{9pt}}}{\tiny\taxid{WT}} & \multicolumn{1}{c}{\tiny\taxid{Muon}} & \multicolumn{1}{c}{\tiny\taxid{MuonT}} & \multicolumn{1}{c@{\hspace{9pt}}}{\tiny\taxid{WT}} & \multicolumn{1}{c}{\tiny\taxid{Muon}} & \multicolumn{1}{c}{\tiny\taxid{MuonT}} & \multicolumn{1}{c@{\hspace{9pt}}}{\tiny\taxid{WT}} & \multicolumn{1}{c}{\tiny\taxid{Muon}} & \multicolumn{1}{c}{\tiny\taxid{MuonT}} & \multicolumn{1}{c@{\hspace{9pt}}}{\tiny\taxid{WT}} & \multicolumn{1}{c}{\tiny\taxid{Muon}} & \multicolumn{1}{c}{\tiny\taxid{MuonT}} & \multicolumn{1}{c@{}}{\tiny\taxid{WT}} \\
\midrule
\multicolumn{23}{@{}l}{\emph{Right-input second moments $C_t = \mathbb{E}[x x^{\Tr}]$ (4 per layer, 128 total)}} \\
\texttt{in:attn} ($q,k,v$) & 512 & $30.0$ & $17.9$ & $18.0$ & $0.63$ & $0.84$ & $0.76$ & $0.1465$ & $0.2031$ & $0.2066$ & $5.8$ & $3.5$ & $3.1$ & $1.4$ & $1.0$ & $0.8$ & $8.2$ & $10.6$ & $10.1$ & $0.041$ & $0.021$ & $0.023$ \\
\texttt{in:o} & 512 & $37.9$ & $20.8$ & $23.4$ & $0.77$ & $0.75$ & $0.77$ & $0.0968$ & $0.2120$ & $0.2001$ & $25.0$ & $5.3$ & $6.5$ & $9.8$ & $1.8$ & $1.8$ & $30.1$ & $22.5$ & $26.5$ & $0.107$ & $0.094$ & $0.108$ \\
\texttt{in:mlp} ($\text{gate},\text{up}$) & 512 & $33.8$ & $21.6$ & $19.5$ & $0.48$ & $0.64$ & $0.62$ & $0.1410$ & $0.1959$ & $0.2024$ & $6.2$ & $3.7$ & $3.2$ & $1.0$ & $0.8$ & $0.7$ & $4.9$ & $6.0$ & $6.7$ & $0.046$ & $0.021$ & $0.025$ \\
\texttt{in:down} (post-SwiGLU) & 2048 & $67.1$ & $29.1$ & $35.6$ & $0.71$ & $2.92$ & $2.33$ & $0.0283$ & $0.0346$ & $0.0359$ & $37.9$ & $47.3$ & $38.8$ & $20.3$ & $43.4$ & $34.7$ & $91.3$ & $79.4$ & $78.7$ & $0.186$ & $0.106$ & $0.087$ \\
\addlinespace[2.5pt] \multicolumn{23}{@{}l}{\emph{Left-cotangent second moments $B_t = \mathbb{E}[\delta \delta^{\Tr}]$ (7 per layer, 224 total)}} \\
\texttt{out:q} & 512 & $38.5$ & $38.5$ & $29.3$ & $1.04$ & $0.97$ & $1.14$ & $0.0615$ & $0.0641$ & $0.0732$ & $27.5$ & $30.0$ & $20.1$ & $16.4$ & $18.0$ & $13.3$ & $64.8$ & $62.4$ & $66.0$ & $0.150$ & $0.140$ & $0.110$ \\
\texttt{out:k} & 512 & $24.4$ & $23.7$ & $13.0$ & $1.45$ & $1.50$ & $1.87$ & $0.0746$ & $0.0774$ & $0.0779$ & $21.7$ & $20.5$ & $18.0$ & $15.5$ & $14.4$ & $14.7$ & $77.6$ & $71.8$ & $78.6$ & $0.178$ & $0.165$ & $0.147$ \\
\texttt{out:v} & 512 & $47.0$ & $50.4$ & $41.5$ & $0.63$ & $0.58$ & $0.57$ & $0.0600$ & $0.0560$ & $0.0749$ & $36.7$ & $41.4$ & $26.4$ & $11.5$ & $12.2$ & $6.8$ & $38.6$ & $31.2$ & $19.9$ & $0.664$ & $0.662$ & $0.586$ \\
\texttt{out:o} (residual) & 512 & $83.3$ & $85.3$ & $78.2$ & $0.09$ & $0.06$ & $0.10$ & $0.0307$ & $0.0301$ & $0.0393$ & $67.1$ & $68.0$ & $55.4$ & $1.0$ & $0.8$ & $0.8$ & $3.6$ & $1.2$ & $2.3$ & $0.164$ & $0.134$ & $0.103$ \\
\texttt{out:down} (residual) & 512 & $84.3$ & $84.6$ & $77.9$ & $0.08$ & $0.07$ & $0.09$ & $0.0315$ & $0.0312$ & $0.0403$ & $65.9$ & $66.3$ & $53.8$ & $0.9$ & $0.8$ & $0.8$ & $2.3$ & $1.2$ & $1.4$ & $0.123$ & $0.107$ & $0.072$ \\
\rowcolor{pagepanel} \textbf{\texttt{out:gate} (pre-SwiGLU)} & 2048 & $88.2$ & $87.1$ & $72.8$ & $0.55$ & $0.59$ & $1.32$ & $\mathbf{0.0055}$ & $\mathbf{0.0051}$ & $\mathbf{0.0048}$ & $94.9$ & $94.9$ & $94.1$ & $\mathbf{80.6}$ & $\mathbf{82.0}$ & $\mathbf{89.5}$ & $\mathbf{98.8}$ & $\mathbf{82.2}$ & $\mathbf{94.9}$ & $0.292$ & $0.240$ & $0.159$ \\
\rowcolor{pagepanel} \textbf{\texttt{out:up} (pre-SwiGLU)} & 2048 & $83.8$ & $82.5$ & $47.3$ & $0.68$ & $0.75$ & $2.83$ & $\mathbf{0.0068}$ & $\mathbf{0.0073}$ & $\mathbf{0.0082}$ & $91.9$ & $91.2$ & $92.4$ & $\mathbf{77.9}$ & $\mathbf{78.1}$ & $\mathbf{90.8}$ & $\mathbf{98.2}$ & $\mathbf{75.9}$ & $\mathbf{91.9}$ & $0.251$ & $0.215$ & $0.130$ \\
\bottomrule
\end{tabular}
}
\end{table}

Table~\ref{tab:motivation} and Figure~\ref{fig:motivation} (top row and
bottom-left) expose three distinct covariance regimes that persist across the
entire $2.68\text{B}$-token training run:

\begin{itemize}
\item \textbf{Width-2048 pre-SwiGLU cotangents (\texttt{out:gate},
  \texttt{out:up}) are coordinate-anisotropic and off-diagonally decorrelated.}
  Across the $32$ layers at update $2{,}500$, $B_{\text{gate}}$ and
  $B_{\text{up}}$ concentrate \motOutGateDiagEMedPct\% and
  \motOutUpDiagEMedPct\% of $\|B\|_F^2$ on their diagonal ($32$-layer means
  \motOutGateDiagEMeanPct\% and \motOutUpDiagEMeanPct\%). Their median
  off-diagonal correlation RMS is \motOutGateCorrRmsMed{} and
  \motOutUpCorrRmsMed{} (bulk Layers $4$--$31$ means:
  \motOutGateBulkCorrRmsMean{} and \motOutUpBulkCorrRmsMean{}; $8$-checkpoint
  trajectory medians: \motOutGateTrajCorrRmsMed{} and
  \motOutUpTrajCorrRmsMed{}), sitting directly on the single-update sampling
  noise floor $N_t^{-1/2} = \motNoiseFloorCorrRms{}$ from $0.13\text{B}$ to
  $2.68\text{B}$ tokens (Figure~\ref{fig:motivation}, bottom-left). Crucially,
  after subtracting $I_d$, \motOutGateDiagAnisoMedPct\% and
  \motOutUpDiagAnisoMedPct\% of $\|\tilde B - I_d\|_F^2$ still resides on the
  diagonal (bulk medians \motOutGateBulkDiagAnisoMedPct\% and
  \motOutUpBulkDiagAnisoMedPct\%; $32$-layer means \motOutGateDiagAnisoMeanPct\%
  and \motOutUpDiagAnisoMeanPct\%), because individual MLP channels exhibit
  large variance spread: mean $\mathrm{CV}_{\diag}$ is \motOutGateDiagCvMean{}
  and \motOutUpDiagCvMean{} (growing from \motOutGateEarlyCvMean{} at update
  $250$ to a peak of \motOutGatePeakCvMean{} at update $3{,}500$), with
  max-to-mean channel ratios averaging $\motOutGateDiagMaxMean\times$ and
  $\motOutUpDiagMaxMean\times$ and reaching $\motOutGateDiagMaxPeak\times$ and
  $\motOutUpDiagMaxPeak\times$ in boundary blocks. Consequently, exact diagonal
  roots $L_{\diag}$ align with dense $L=A^{-1/2}$ at median cosine
  \motOutGateCosLMed{} and \motOutUpCosLMed{} and explain
  \motOutGateExpVarLMedPct\% and \motOutUpExpVarLMedPct\% of $\|L-I_d\|_F^2$
  (means \motOutGateExpVarLMeanPct\% and \motOutUpExpVarLMeanPct\%; trajectory
  averages \motOutGateTrajExpVarLMedPct\% and \motOutUpTrajExpVarLMedPct\%),
  whereas omitting preconditioning ($L=I_d$) incurs more than double the
  relative error (\motOutGateDropErrLMean{} and \motOutUpDropErrLMean{} versus
  \motOutGateRelErrLMean{} and \motOutUpRelErrLMean{}).

\item \textbf{Width-512 residual-stream cotangents (\texttt{out:o},
  \texttt{out:down}) are $>99\%$ rotational.} Because \texttt{o} and
  \texttt{down} add directly into the shared residual stream, their per-channel
  cotangent variances are nearly uniform ($\mathrm{CV}_{\diag} = \motOutODiagCvMed{}$
  and $\motOutDownDiagCvMed{}$ median, $\motOutODiagCvMean{}$ and
  $\motOutDownDiagCvMean{}$ mean). Although raw $\mathrm{DiagE}$ appears high
  (\motOutODiagEMedPct\% and \motOutDownDiagEMedPct\%) due to a large isotropic
  floor ($\mathrm{erank}/d = \motOutOErankMedPct\%$ and
  $\motOutDownErankMedPct\%$), subtracting $I_d$ collapses
  $\mathrm{DiagAniso}$ to \motOutODiagAnisoMedPct\% and
  \motOutDownDiagAnisoMedPct\% across the entire trajectory
  (Figure~\ref{fig:motivation}, top-right), and $L_{\diag}$ explains only
  \motOutOExpVarLMedPct\% and \motOutDownExpVarLMedPct\% of $\|L-I_d\|_F^2$
  (means \motOutOExpVarLMeanPct\% and \motOutDownExpVarLMeanPct\%). Attention
  cotangents sit between these extremes: \texttt{out:q} and \texttt{out:k} carry
  substantial head/RoPE coordinate spread ($R^2 = \motOutQExpVarLMedPct\%$ and
  $\motOutKExpVarLMedPct\%$), while \texttt{out:v} is predominantly rotational
  ($R^2 = \motOutVExpVarLMedPct\%$).

\item \textbf{Right-input moments ($C_t$) are strongly coupled off-diagonal,
  including at width 2048.} Preceded by RMSNorm or attention mixing, width-$512$
  inputs (\texttt{in:attn}, \texttt{in:o}, \texttt{in:mlp}) have low effective
  rank ($\motInAttnErankMedPct\%$--$\motInMlpErankMedPct\%$), strong
  cross-channel correlations ($\rho_{\mathrm{rms}} = \motInMlpCorrRmsMed{}$--$\motInOCorrRmsMed{}$,
  $30\times$--$40\times$ above \texttt{out:gate}), and only
  $\motInMlpDiagAnisoMedPct\%$--$\motInODiagAnisoMedPct\%$ centered diagonal
  anisotropy ($98\%$--$99\%$ off-diagonal). Even at width $2048$, post-SwiGLU
  \texttt{in:down} behaves quite differently from pre-SwiGLU \texttt{out:gate}
  and \texttt{out:up}: despite massive neuron scale spread (mean
  $\mathrm{CV}_{\diag}$ growing from \motInDownEarlyCvMean{} to
  \motInDownFinalCvMean{}, giving high $R^2(A^{-1/2}) = \motInDownExpVarLMedPct\%$
  purely from coordinate variance heterogeneity), its $\mathrm{erank}/d$ is only
  \motInDownErankMedPct\%, $\rho_{\mathrm{rms}}$ is \motInDownCorrRmsMed{}
  ($6.8\times$ \texttt{out:gate}'s RMS, or $46\times$ higher correlation energy
  per coordinate), and $56.6\%$ of $\|\tilde C_{\text{down}}-I_d\|_F^2$ lies off
  the diagonal ($\mathrm{DiagAniso} = \motInDownDiagAnisoMedPct\%$).
\end{itemize}

\begin{finding}[Hadamard decorrelation explains the pre- versus post-SwiGLU split at width 2048]
In a SwiGLU MLP $y = W_{\text{down}}(\operatorname{SiLU}(g)\odot u)$ with
$g=W_{\text{gate}}x$ and $u=W_{\text{up}}x$, the pre-activation cotangents are
elementwise Hadamard products
$\delta_{\text{up}} = (\delta_{\text{down}}W_{\text{down}})\odot \operatorname{SiLU}(g)$
and $\delta_{\text{gate}} = (\delta_{\text{down}}W_{\text{down}})\odot (u\odot \operatorname{SiLU}'(g))$
between a zero-mean backward residual vector and a forward gate vector. Under
Schur multiplication of weakly dependent covariances
($\Sigma_{ij} \approx \Sigma^{\text{bwd}}_{ij}\Sigma^{\text{fwd}}_{ij}$), small
off-diagonal correlations multiply toward zero ($|\rho_{ij}^{\text{bwd}}\rho_{ij}^{\text{fwd}}| < |\rho_{ij}|$)
down to the $N_t^{-1/2}$ sampling floor, whereas diagonal variances multiply
together and compound neuron scale anisotropy. By contrast, the forward
post-SwiGLU input $x_{\text{down}} = \operatorname{SiLU}(g)\odot u$ has strictly
positive channel means from one-sided $\operatorname{SiLU}$, adding a dense
outer-product mean term $\mu\mu^{\Tr}$ to $\mathbb E[x_{\text{down}}x_{\text{down}}^{\Tr}]$
that keeps $C_{\text{down}}$ strongly correlated off-diagonal --- a sharp,
falsifiable prediction tested directly by symmetric width-2048 diagonalization
(\taxid{WT-Sym}, Section~\ref{sec:sym-ablation}).
\end{finding}

\begin{figure}[t]
\centering
\resizebox{\textwidth}{!}{%
\begin{tikzpicture}
\begin{groupplot}[
  paperplot,
  group style={group size=2 by 2, horizontal sep=1.75cm, vertical sep=1.70cm},
  width=7.4cm, height=4.9cm,
]
\nextgroupplot[
  ybar=0.8pt, bar width=2.4pt,
  ylabel={Share / fidelity (\%)},
  ymin=0, ymax=108,
  xmin=-0.6, xmax=10.6,
  xtick={0,1,2,3,4,5,6,7,8,9,10},
  xticklabels={\texttt{in:at},\texttt{in:o},\texttt{in:mlp},\texttt{in:dn},\texttt{q},\texttt{k},\texttt{v},\texttt{o},\texttt{down},\texttt{gate},\texttt{up}},
  xticklabel style={rotate=36, anchor=east, font=\tiny},
  title={(a) Diagonal energy, anisotropy \& root $R^2$},
  legend columns=3,
  legend style={at={(0.5,0.98)}, anchor=north, font=\tiny, column sep=3pt},
]
\addplot[fill=famSlate!45, draw=famSlate!80]
  table[x=idx, y=diag_energy_med_pct] {figs/motivation_bars.dat};
\addlegendentry{$\mathrm{DiagE}$}
\addplot[fill=famNull!78, draw=famNull!95]
  table[x=idx, y=diag_aniso_med_pct] {figs/motivation_bars.dat};
\addlegendentry{$\mathrm{DiagAniso}$}
\addplot[fill=famWT!78, draw=famWT!95]
  table[x=idx, y=exp_var_L_med_pct] {figs/motivation_bars.dat};
\addlegendentry{$R^2(A^{-1/2})$}

\nextgroupplot[
  xlabel={Tokens processed (billions)},
  ylabel={$\mathrm{DiagAniso}(\tilde M)$ (\%)},
  xmin=0.05, xmax=2.75, ymin=-3, ymax=122,
  ytick={0,25,50,75,100},
  title={(b) Anisotropic diagonal share over training},
  legend columns=3,
  legend style={at={(0.5,0.98)}, anchor=north, font=\tiny, column sep=2.5pt},
  every axis plot/.append style={line width=0.9pt, mark size=1.4pt},
]
\addplot[color=famOD, mark=*] table[x=tokens, y=da_gate] {figs/motivation_traj.dat};
\addlegendentry{\texttt{out:gate}}
\addplot[color=famOD!70, densely dashed, mark=o] table[x=tokens, y=da_up] {figs/motivation_traj.dat};
\addlegendentry{\texttt{out:up}}
\addplot[color=famNull, mark=triangle*] table[x=tokens, y=da_in_down] {figs/motivation_traj.dat};
\addlegendentry{\texttt{in:down}}
\addplot[color=famWT, mark=square*] table[x=tokens, y=da_out_qk] {figs/motivation_traj.dat};
\addlegendentry{\texttt{out:q,k}}
\addplot[color=famRef!75, densely dotted, mark=diamond*] table[x=tokens, y=da_in_d512] {figs/motivation_traj.dat};
\addlegendentry{\texttt{in:512}}
\addplot[color=famAbl, mark=x] table[x=tokens, y=da_out_odown] {figs/motivation_traj.dat};
\addlegendentry{\texttt{out:o,down}}

\nextgroupplot[
  xlabel={Tokens processed (billions)},
  ylabel={Off-diagonal $\rho_{\mathrm{rms}}$ (log scale)},
  ymode=log,
  xmin=0.05, xmax=2.75, ymin=0.0022, ymax=0.45,
  title={(c) Correlation RMS vs.\ $N_t^{-1/2}$ floor},
  legend columns=3,
  legend style={at={(0.5,0.56)}, anchor=center, font=\tiny, column sep=2pt},
  every axis plot/.append style={line width=0.9pt, mark size=1.4pt},
]
\addplot[color=famRef!75, densely dotted, mark=diamond*] table[x=tokens, y=rho_in_d512] {figs/motivation_traj.dat};
\addlegendentry{\texttt{in:512}}
\addplot[color=famWT, mark=square*] table[x=tokens, y=rho_out_qk] {figs/motivation_traj.dat};
\addlegendentry{\texttt{out:q,k}}
\addplot[color=famNull, mark=triangle*] table[x=tokens, y=rho_in_down] {figs/motivation_traj.dat};
\addlegendentry{\texttt{in:down}}
\addplot[color=famAbl, mark=x] table[x=tokens, y=rho_out_odown] {figs/motivation_traj.dat};
\addlegendentry{\texttt{out:o,down}}
\addplot[color=famOD!70, densely dashed, mark=o] table[x=tokens, y=rho_up] {figs/motivation_traj.dat};
\addlegendentry{\texttt{out:up}}
\addplot[color=famOD, mark=*] table[x=tokens, y=rho_gate] {figs/motivation_traj.dat};
\addlegendentry{\texttt{out:gate}}
\draw[checkgreen!85!black, dashed, line width=0.8pt] (axis cs:0.05,0.005226) -- (axis cs:2.75,0.005226);
\node[font=\tiny, text=checkgreen!80!black, anchor=north west]
  at (axis cs:0.08,0.0049) {$N_t^{-1/2}=0.0052$};

\nextgroupplot[
  xlabel={Tokens processed (billions)},
  ylabel={8-GPU exact \texttt{eigh} fallbacks ($\times 10^3$)},
  xmin=0.05, xmax=2.75, ymin=-4, ymax=118,
  title={(d) Cumulative exact fallbacks ($B_t$ vs.\ $C_t$)},
  legend style={at={(0.03,0.97)}, anchor=north west, font=\tiny},
  every axis plot/.append style={line width=1.05pt, mark size=1.6pt},
]
\addplot[color=famAbl, mark=*] table[x=tokens, y=out_fb_k] {figs/motivation_traj.dat};
\addlegendentry{Left cotangents $B_t$ ($\motFbOutTotal$)}
\addplot[color=famWT, mark=square*] table[x=tokens, y=in_fb_k] {figs/motivation_traj.dat};
\addlegendentry{Right inputs $C_t$ ($\motFbInTotal$)}
\node[font=\tiny, text=pageink!85, anchor=west, align=left]
  at (axis cs:0.15,66) {After $0.13\text{B}$ tok.:\\$\motFbOutPostWarmup$ vs.\ $\motFbInPostWarmup$ ($\motFbRatioPostWarmup\times$)};
\end{groupplot}
\end{tikzpicture}%
}
\caption{Full-trajectory covariance and temporal-stability audit of baseline
\taxid{MuonT} (\nullall{}, eight checkpoints from $0.13\text{B}$ to $2.68\text{B}$
tokens across all $32$ layers). \textbf{(a)} Subtracting $I_d$ separates
genuinely diagonal moments (\texttt{out:gate}, \texttt{out:up}) from isotropic
but rotation-dominated residual moments (\texttt{out:o}, \texttt{out:down}),
while showing why \texttt{in:down} still retains $56.6\%$ off-diagonal
anisotropy despite high scale spread. \textbf{(b)} The architectural split in
$\mathrm{DiagAniso}(\tilde M)$ holds throughout training. \textbf{(c)}
Pre-SwiGLU \texttt{out:gate} and \texttt{out:up} hug the single-step sampling
floor $N_t^{-1/2}=0.0052$ across the entire run, whereas width-$512$ inputs sit
$35\times$ higher and \texttt{in:down} sits $6.8\times$ higher. \textbf{(d)}
Rapid cotangent subspace rotation triggers $\motFbOutTotal$ exact
\texttt{cuSOLVER eigh} fallbacks across the eight GPUs versus $\motFbInTotal$ on
inputs (only $\motFbInPostWarmup$ input fallbacks after $0.13\text{B}$ tokens).}
\label{fig:motivation}
\end{figure}

\subsection{Temporal drift of cotangent eigenspaces}
\label{sec:mot-drift}

Spatial covariance structure is only half of the design problem; the second half
is temporal stability. Preconditioned \muont{} avoids per-update
eigendecompositions by expressing each current moment $\tilde M_t$ in a cached
orthonormal eigenbasis $X_{t-\tau}$ refreshed asynchronously on host worker
threads every four updates, refining $(X_{t-\tau}^{\Tr}\tilde M_t X_{t-\tau} + \epsilon I_d)^{-1/4}$
via two coupled fixed-point iterations (up to three repair rounds when the
estimated error exceeds tolerance, falling back to synchronous GPU
\texttt{cuSOLVER eigh} when repair fails; Section~\ref{sec:method}). That scheme
was engineered for forward activations $x_t$, which are stabilized by RMSNorm.
Backward cotangents $\delta_t$ enjoy no such normalization and depend on
changing downstream weights and token-level loss residuals.

To quantify how fast the two sides rotate along the \muont{} trajectory, we
measure the top-$k$ Grassmann subspace drift between consecutive checkpoints
$\Delta t = 250$ updates apart,
\begin{equation}
\label{eq:grassmann}
\mathrm{Drift}_k(t_1, t_2)
= 1 - \frac{1}{k}\bigl\|V_{1:k}(t_1)^{\Tr} V_{1:k}(t_2)\bigr\|_F^2 \in [0, 1],
\end{equation}
where $V_{1:k}(t)\in\mathbb R^{d\times k}$ spans the $k$ principal eigenvectors
of $\tilde M_t$. Across all $32$ layers between updates $2{,}250$ and $2{,}500$
(last column of Table~\ref{tab:motivation}), RMSNormed input subspaces barely
move: $\mathrm{Drift}_{16}$ is \motInAttnDriftTopSixteenMean{} on
\texttt{in:attn} and \motInMlpDriftTopSixteenMean{} on \texttt{in:mlp}
($\mathrm{Drift}_{32} = \motInAttnDriftTopThirtyTwoMean{}$ and
$\motInMlpDriftTopThirtyTwoMean{}$), and \motInODriftTopSixteenMean{}--\motInDownDriftTopSixteenMean{}
on \texttt{in:o} and \texttt{in:down}. Left-cotangent subspaces rotate
$5\times$ to $\motDriftRatioOutVVsInAttn\times$ faster:
$\mathrm{Drift}_{16}$ reaches \motOutQDriftTopSixteenMean{} and
\motOutKDriftTopSixteenMean{} on \texttt{out:q} and \texttt{out:k},
\motOutUpDriftTopSixteenMean{} and \motOutGateDriftTopSixteenMean{} on
\texttt{out:up} and \texttt{out:gate} ($\motDriftRatioOutUpVsInAttn\times$ and
$\motDriftRatioOutGateVsInAttn\times$ \texttt{in:attn}), and
\motOutVDriftTopSixteenMean{} on \texttt{out:v} ($\mathrm{Drift}_{32} = \motOutVDriftTopThirtyTwoMean{}$).
The same gap holds in early training (updates $250\to 500$:
\motOutGateEarlyDriftTopSixteen{} on \texttt{out:gate} versus
\motInAttnEarlyDriftTopSixteen{} on \texttt{in:attn}) and late training (updates
$4{,}250\to 4{,}500$: \motOutVLateDriftTopSixteen{} on \texttt{out:v} versus
\motInAttnLateDriftTopSixteen{} on \texttt{in:attn}). Even over the $1$--$4$
update age of the saved cached bases $X_{t-\tau}$ inside the update-$2{,}500$
checkpoint, short-horizon top-$16$ drift is $24\times$--$36\times$ larger on
\texttt{out:v} and \texttt{out:gate} (\motOutVShortDriftTopSixteen{} and
\motOutGateShortDriftTopSixteen{}) than on \texttt{in:attn} and \texttt{in:mlp}
(\motInAttnShortDriftTopSixteen{} and \motInMlpShortDriftTopSixteen{}).

Evaluating the two-iteration fixed-point root solver directly on the cached
eigenbases $X_{t-\tau}^{\Tr}\tilde M_t X_{t-\tau}$ saved at update $2{,}500$
shows the immediate numerical consequence. On the input side, the median ratio of
estimated root error to acceptance tolerance is
\motInAttnCachedEstOverTolMed{}--\motInOCachedEstOverTolMed{}, comfortably below
$1.0$. On the output side, \texttt{out:k} sits at median ratio
\motOutKCachedEstOverTolMed{} and \texttt{out:v} at
\motOutVCachedEstOverTolMed{}, while each rank's two host worker threads must
simultaneously service $28$ output moments including eight width-$2048$
eigendecompositions. Figure~\ref{fig:motivation}(d) shows the resulting
run-wide divergence across all eight GPUs: over $5{,}120$ updates, dense input
moments incur $\motExtraInTotal{}$ repair rounds and only $\motFbInTotal{}$
exact \texttt{eigh} fallbacks ($\motFbInPctStatStep\%$ of evaluations), whereas
dense output moments trigger $\motExtraOutTotal{}$ repair rounds and
$\motFbOutTotal{}$ exact fallbacks ($\motFbOutPerUpdate{}$ per update globally,
$\motFbOutPctStatStep\%$ of all $1.14\times 10^6$ output evaluations ---
$\motFbRatioTotal\times$ more fallbacks, or $\motFbRatioPerStat\times$ per
statistic, and $\motFbRankZeroOut{}$ versus $\motFbRankZeroIn{}$ on Rank~0).
After the initial $250$-update ($0.13\text{B}$-token) transient, all eight GPUs
record just $\motFbInPostWarmup{}$ input fallbacks over the remaining $4{,}870$
updates against $\motFbOutPostWarmup{}$ output fallbacks --- a
$\motFbRatioPostWarmup\times$ raw ($\motFbRatioPostWarmupPerStat\times$
per-statistic) gap. Because backward cotangents become available only at the end
of the last microbatch's backward pass, executing these output repairs and
\texttt{cuSOLVER eigh} fallbacks synchronously stalls the optimizer step
directly on the critical path.

\subsection{Three design consequences}
\label{sec:mot-consequences}

Taken together, the spatial decomposition in Table~\ref{tab:motivation} and the
temporal drift in Figure~\ref{fig:motivation}(d) dictate three concrete design
choices before any candidate optimizer is trained:

\begin{enumerate}
\item \textbf{Keep width-512 input whitening dense, and treat width-2048
  \texttt{in:down} as a distinct quality--speed trade-off.} Between $98\%$
  and $99\%$ of width-$512$ input anisotropy --- and $56.6\%$ of post-SwiGLU
  \texttt{in:down}'s --- lies off the coordinate diagonal. Dropping or
  diagonalizing width-$512$ inputs forfeits the primary domain conditioning of
  the finite polynomial map $\Nten$; diagonalizing \texttt{in:down} alongside
  \texttt{out:gate,up} (\taxid{WT-Sym}) eliminates all width-$2048$ dense Gram
  and eigendecomposition work on \emph{both} sides, but unlike pre-SwiGLU
  outputs it discards real off-diagonal $\mu\mu^{\Tr}$ structure.
\item \textbf{Threshold output richness by width (\taxid{WT-MLP} / \mlpdiag{}),
  with pure diagonals (\taxid{OD-All} / \alldiag{}) as the
  zero-eigendecomposition limit.} The two width-$2048$ pre-SwiGLU cotangents
  (\texttt{out:gate}, \texttt{out:up}) account for $96.2\%$ of output $O(d^3)$
  eigendecomposition work and $86.5\%$ of $O(d^2)$ Gram work, yet Hadamard
  decorrelation pins their off-diagonal correlations to the $N_t^{-1/2}$
  sampling floor while leaving strong per-neuron scale spread. Replacing their
  dense spectral roots with exact, age-zero elementwise diagonals captures
  $\approx 76\%$--$82\%$ of $\|L-I_d\|_F^2$ on baseline trajectories without
  fitting $2.1\times 10^6$ noisy off-diagonals per matrix or clogging host
  worker threads. Conversely, width-$512$ residual-stream cotangents
  (\texttt{out:o}, \texttt{out:down}) are $>99\%$ rotational ($R^2 =
  \motOutOExpVarLMedPct\%$), motivating dense spectral roots at width $512$
  (\taxid{WT-MLP}), while diagonalizing all seven outputs (\taxid{OD-All})
  defines an eigendecomposition-free speed bound via \texttt{gate},
  \texttt{up}, \texttt{q}, and \texttt{k}.
\item \textbf{Lag output factors by one update on an asynchronous side stream
  (\taxid{WT-MLP-Lag}, \taxid{WT-Sym-Lag}).} Diagonalizing width-$2048$ outputs
  eliminates the largest host bottlenecks, but fast subspace drift on width-$512$
  cotangents (\texttt{out:v}, \texttt{out:k}) still triggers frequent iterative
  repairs and exact \texttt{eigh} fallbacks ($\approx 20{,}000$--$22{,}000$ on
  Rank~0 per run). Applying factors $(L^{(t-1)},\hat L^{(t-1)})$ built from the
  previous update's cotangents --- a staleness of one-sixteenth of the
  $\beta=0.938$ EMA window --- moves output Gram reduction, root refinement, and
  fallback resolution off the synchronous post-backward critical path and
  overlaps them with the next update's forward and backward pass.
\end{enumerate}
